\subsection{学习表示与利用表示结构}
OPQ通过学习空间分解和码本降低失真\cite{opq}，其旋转为数据优化；
JHQ随机平方旋转支持廉价解析primary模型。取舍涉及拟合成本、重建保真和评分结构，
更好的通用学习码本未必保留16+16精确分解所需的内部Cartesian乘积。

ScaNN的各向异性量化根据最大内积检索的重要性加权误差，说明重建MSE不是唯一有用
目标\cite{scann}。其目标与实现环境不同于本文平方L2层次流水线。
可比的是损失选择：面向排序的表示仍须检验实际近邻保留及执行成本。
第8章停止的排序研究在JHQ内部提供了例子。

\subsection{小查找表已有充分先行工作}
PQ Fast Scan把依赖cache的表改造成适于CPU SIMD寄存器的表\cite{fastscan}；
Quicker ADC推广寄存器查找，引入不规则量化器与split tables\cite{quickeradc}。
这些工作说明表形状要与执行指令共同设计。本文的范围是继承Cartesian坐标结构所支持
的精确分解及其GPU实测；CPU SIMD结果提供概念先例，不是本GPU的吞吐。
FP32结合顺序和cutoff并列仍需独立数值检查。

RaBitQ也使用随机化量化与紧凑符号信息，具有特定距离估计器及理论误差分析
\cite{rabitq}；多bit扩展以存储换估计精度\cite{rabitq_multi}。
其保证属于相应构造和假设，不能因双方都用旋转与二值组件就转移给JHQ。
JHQ主符号选survivor，显式存储的残差提供下一精度层；质量同时依赖截断与重建。

GPU IVF-RaBitQ把该族整合成GPU建库与搜索，包括GPU距离计算和融合搜索
\cite{ivfrabitq}，因此是直接相关的现代基线，不是“缺少GPU量化器”的背景证据。
其融合执行说明紧凑编码还需要有竞争力的访问和选择流水线。
E1纳入多个bit预算，其高recall结果限制JHQ的宽泛优势主张。

\subsection{精排工作量与自适应停止}
候选扩张和重排已经存在于Faiss的\texttt{k\_factor}及cuVS refinement控制
\cite{faiss_refine,cuvs_refinement}。Faiss autotuning可用GT标签搜索recall--时间取舍
\cite{faiss_autotune}；学习式自适应终止预测足够搜索量\cite{li_termination}，
DARTH围绕声明式recall目标组织早停\cite{darth}。
这些都是均匀工作量之外的相关替代。第6章的具体策略用大预算参考的无标签一致性，
为整个工作负载选择一个survivor倍率，省去预测器训练，却继承参考的质量局限。
这是策略范围，不表示自适应搜索或精排预算是新概念。
